% Compile with XeLaTeX. Submission specification: 12 pt, double-spaced, 4--6 pages.
% Revision uses the attached proposal, coauthor review, and supplied project discussion.
% DeepSeek share could not be retrieved; no content is attributed to that discussion.
\documentclass[12pt,a4paper]{article}
\usepackage[margin=1in]{geometry}
\usepackage{fontspec}
\IfFontExistsTF{Times New Roman}
  {\setmainfont{Times New Roman}}
  {\setmainfont[
    Path=/usr/share/fonts/opentype/urw-base35/,
    Extension=.otf,
    UprightFont=NimbusRoman-Regular,
    BoldFont=NimbusRoman-Bold,
    ItalicFont=NimbusRoman-Italic,
    BoldItalicFont=NimbusRoman-BoldItalic
  ]{NimbusRoman-Regular}}
\usepackage{amsmath,amssymb,setspace,enumitem,titlesec}
\usepackage[hidelinks]{hyperref}
\doublespacing
\titleformat{\section}{\normalfont\normalsize\bfseries}{\thesection.}{0.5em}{}
\titlespacing*{\section}{0pt}{6pt}{2pt}
\setlength{\parindent}{0pt}
\setlength{\parskip}{3pt}
\setlength{\emergencystretch}{2em}
\setlist{nosep,leftmargin=*}
\hypersetup{pdftitle={Symmetry-Adapted Neural Surrogates and Certified Lie Reduction for Saint--Venant--Exner Flows}}
\begin{document}
\begin{center}
\textbf{Symmetry-Adapted Neural Surrogates and Certified Lie Reduction}\\
\textbf{for Saint--Venant--Exner Flows}\\
\textit{A parameter-conditioned series--neural study}\\
Carl \quad WKU MATH 4890\\
Proposed supervisor: Dr. Puneet Rana
\end{center}
\section{Research problem and objectives}
The Saint--Venant--Exner (SVE) equations describe shallow-water motion coupled to sediment-bed evolution. This project asks whether a verified Lie reduction can support a \textbf{parameter-conditioned neural surrogate} for a smooth family of water--sediment flows, while retaining computable bounds on depth, velocity, and bed elevation.

The physical model has one spatial coordinate and time: it is a $(1+1)$-dimensional nonlinear PDE system with three fields. Lie reduction selects a special family and removes the spatial independent variable; it does not solve arbitrary initial data. The reduced problem remains a coupled nonlinear ODE system evolving in time.

The cubic SVE invariant family is already studied in \cite{sve2022}, and Taylor polynomials with neural remainders are established in \cite{taylor}. The proposed contribution is a coupled Lie--ML--analysis study: hard symmetry-adapted coordinates for a neural solution map, a local series that fixes its known behavior, and model-specific error transfer to physical observables. A generic Gr\"onwall inequality or a new combination of existing solvers is not itself a novelty claim.

The reduction, certification protocol, and neural surrogate are coequal deliverables. The reduction supplies the coordinates and equations that make the learning problem precise; certification determines when the surrogate has a physical meaning. Classical integration remains the required baseline, not the intended endpoint.
\newpage
\section{Governing model and exact Lie reduction}
In nondimensional variables, let $h>0$ denote depth, $u$ velocity, and $B$ bed elevation. Fix $g=1$ and $\kappa=0.1$ for the initial experiments:
\begin{equation}\begin{aligned}
h_t+(hu)_x&=0,\\
u_t+uu_x+g(h+B)_x&=-d(t;\varepsilon)u,\\
B_t+\kappa(u^3)_x&=0.
\end{aligned}\label{pde}\end{equation}
Here $\kappa=A_g/(1-p)$ includes sediment porosity $p$ and transport coefficient $A_g$. The baseline has $d=0$. An optional prescribed-drag extension is
\[
d(t;\varepsilon)=\varepsilon[1+\tfrac12\sin(2\pi(t-1))],\quad 0\leq\varepsilon\leq0.05.
\]
This is an idealized continuum model, not a molecular or chemical model.

The generator $X=x\partial_x+u\partial_u+2h\partial_h+2B\partial_B$ scales $(x,u,h,B)$ by $(\lambda,\lambda,\lambda^2,\lambda^2)$ and fixes $t$. The mass and bed balances scale by $\lambda^2$, and momentum by $\lambda$, including drag. The invariant conditions $xu_x=u$, $xh_x=2h$, $xB_x=2B$ integrate to
\[
u=xv(t),\qquad h=x^2H(t),\qquad B=x^2C(t).
\]
For example, $(hu)_x=3x^2Hv$ and $(u^3)_x=3x^2v^3$. Substitution gives $Y'=f(t,Y;\mu)$ for $Y=(v,H,C)^T$:
\begin{equation}
v'=-v^2-2g(H+C)-dv,\quad H'=-3vH,\quad C'=-3\kappa v^3.\label{ode}
\end{equation}
Construct the family on $x>0$, then restrict to $1\leq x\leq2$, $1\leq t\leq1.5$. The fixed window is not scaling-invariant; its boundary traces must come from the reconstructed solution.
\newpage
\section{Verification and the proposed error result}
The Lagrange-variable paper \cite{sve} motivates a separate coordinate audit. For $x=\phi(t,s)$, $\phi_t=u(t,\phi)$, $\phi_s=1/h(t,\phi)$, and $\widetilde B=B(t,\phi)$, the chain rule requires
\[
\widetilde B_t-\frac{\phi_t}{\phi_s}\widetilde B_s
 +\frac{\kappa}{\phi_s}\partial_s(\phi_t^3)=0.
\]
The transport term cannot be dropped for this definition of bed height. The proposed Eulerian reduction avoids relying on an unreconciled scalar master equation.

\textbf{Target proposition (conditional error and admissibility bound).}
Let $\widehat Y\in C^1([1,T])$ approximate \eqref{ode}, with defect $R=\widehat Y'-f(t,\widehat Y;\mu)$. Suppose a compact convex region $K$ contains the approximation and a validated enclosure of the exact solution; let
\[
\|R\|_\infty\leq\delta,\quad
\|Y(1)-\widehat Y(1)\|_\infty\leq e_0,\quad
\sup_{[1,T]\times K}\|D_Yf\|_\infty\leq L.
\]
Then, for $z=t-1$,
\[
\|Y(t)-\widehat Y(t)\|_\infty\leq E(t)
 :=e^{Lz}e_0+\delta\frac{e^{Lz}-1}{L},
\]
with $E=e_0+\delta z$ if $L=0$. On $1\leq x\leq2$, this gives
\[
|u-\widehat u|\leq2E,\qquad
|h-\widehat h|\leq4E,\qquad |B-\widehat B|\leq4E.
\]
If $\widehat H(t)-E(t)\geq c>0$, both depths are positive. The proof uses the integral error equation and Gr\"onwall's inequality. The research work is to justify $K,L,\delta,c$ computationally, rather than assume them from a plot. A bound that fails does not prove blow-up or dry-bed formation.
\newpage
\section{Series construction and computational certification}
Use the exact control
\[
Y_*(t)=\left(\frac{2}{3t},\frac{1}{15t^2},\frac{2}{45t^2}\right)^T,\qquad d=0,
\]
verified by substitution in all three balances. For the substantive case, perturb the exact initial profile while retaining $d=0$. A refined adaptive ODE solution supplies a numerical comparison, not a proof. The neural task will then represent a compact family obtained by varying all three reduced initial profiles and the drag amplitude. The initial certification target is uniform physical-field error below $10^{-4}$ on the stated window; this is a proposed tolerance, not an achieved result.

For $z=t-1$, calculate coefficients from the ODE, beginning with
\[
P_2(z)=Y_0+zf_0+\frac{z^2}{2}(f_t+D_Yf\,f)_0.
\]
The subscript $0$ means evaluation at $(1,Y_0;\mu)$. The $f_t$ term accounts for time-dependent drag. Higher coefficients follow by differentiation or coefficient matching. Analyticity gives a local Taylor representation; a useful interval and remainder bound still require proof.

A practical Jacobian bound on $K$ is obtained from
\[
L=\sup_{[1,T]\times K}\max\{,|2v+d|+4g,\ 3|H|+3|v|,\ 9\kappa v^2\,\}.
\]
Use outward-rounded interval bounds for derivatives and residuals on successive time cells. Alternatively, a grid residual bound $\delta_g$, a justified derivative bound $\|R'\|_\infty\leq M$, and a grid covering radius $q$ give $\delta\leq\delta_g+Mq$. Training samples alone cannot provide these bounds.

A validated ODE enclosure or a first-exit argument must ensure that the error tube stays within $K$. Subdivide the interval when necessary and propagate endpoint error. Since $H(t)=H(1)\exp(-3\int_1^t v(s)\,ds)$, regular exact trajectories preserve positive depth. Certification must also establish a quantitative margin for their approximations.
\newpage
\section{Symmetry-adapted neural surrogate, outcomes, and work plan}
The neural surrogate is a central deliverable. It approximates the series remainder for the entire reduced parameter family:
\[
Y_\theta=P_N(z;\mu)+z^{N+1}D\mathcal N_\theta(z;\mu),\quad
D=\operatorname{diag}(2/3,1/15,2/45).
\]
The network learns the three reduced profiles across the parameter family; it does not discover the supplied symmetry. The prefactor enforces the initial value and first $N$ derivatives. Its input is the quotient coordinate $z=t-1$ together with the reduced initial state and drag amplitude, rather than the original spatial coordinate $x$. Train a small smooth network using
\[
\mathcal L=\frac1{N_c}\sum_i
\left\|D^{-1}\big[Y_\theta'-f(t_i,Y_\theta;\mu_i)\big]\right\|_2^2.
\]
Compare series alone, adaptive ODE integration, $Y_0+zD\mathcal N_\theta$, and the series--neural trial \cite{taylor,lagaris}. Match accuracy targets and report setup, training, evaluation, and certification costs. The question is whether a trained surrogate earns its setup cost across many held-out parameter queries; a single reduced ODE solve is expected to remain the stronger baseline. The parameter experiment will use $\mu=(v(1),H(1),C(1),\varepsilon)$ in a compact box around the exact control, with classical interpolation as an additional baseline. Hold out complete parameter cases; pointwise certificates do not certify the entire parameter domain.

\textbf{Expected outcomes.} Deliver a checked reduction, benchmark trajectories, a model-specific error/admissibility proposition, and reproducible comparisons. Report depth, velocity, bed error, positivity margins, and bed change $\Delta B=x^2[C(T)-C(1)]$. The main ML result is a held-out parameter study, measured in profile and reconstructed-field error, residual, positivity margin, and amortized cost. Neural improvement is a question to measure; an informative negative result remains useful.

\textbf{Twelve-week plan.} Weeks 1--2: model, symmetry, reconstruction, and exact control. Weeks 3--4: series and parameter-family design. Weeks 5--7: neural surrogate and ablations. Weeks 8--9: validated enclosures and physical error transfer. Weeks 10--11: held-out parameter tests and cost analysis. Week 12: reproducibility and writing. If rigorous bounds remain incomplete, label the computations as verification and the theorem as conditional.

\textbf{Later direction.} The cubic exponent is special: for a bedload law $u^m$, this spatial scaling exactly balances the Exner equation only at $m=3$. Thus $m=3+\delta$ is an approximate-symmetry problem, not merely another parameter in the reduced ODE. A later study can factor out the cubic solution and learn its symmetry-breaking correction in $(t,\log x,\delta)$. Study certified continuation and perturbations outside the invariant family. The supplied Navier--Stokes manuscript \cite{ns} motivates distinguishing local analytic series from controlled asymptotic corrections. Its stated forced-flow result is not a theorem about this SVE model; full-PDE stability, shocks, and singularity proofs remain separate projects.
\section*{References}
\begin{thebibliography}{9}
\raggedright
\bibitem{sve2022} A. Paliathanasis. ``Lie Symmetry Analysis of the One-Dimensional Saint-Venant--Exner Model.'' \emph{Symmetry} 14, 1679 (2022). \href{https://doi.org/10.3390/sym14081679}{doi:10.3390/sym14081679}.
\bibitem{sve} A. Paliathanasis, G. Leon, and P. G. L. Leach. ``Lie Symmetries and Similarity Solutions for a Shallow-Water Model with Bed Elevation in Lagrange Variables.'' \emph{Mathematics} 14, 433 (2026). \href{https://doi.org/10.3390/math14030433}{doi:10.3390/math14030433}.
\bibitem{taylor} C. K. Nagesh, S. Sankaranarayanan, R. Kaur, T. Sahai, and S. Jha. ``Taylor-Model Physics-Informed Neural Networks (PINNs) for Ordinary Differential Equations.'' \emph{PMLR} 288, 621--642 (2025). \href{https://proceedings.mlr.press/v288/nagesh25a.html}{Open-access proceedings}.
\bibitem{lagaris} I. E. Lagaris, A. Likas, and D. I. Fotiadis. ``Artificial Neural Networks for Solving Ordinary and Partial Differential Equations.'' \emph{IEEE Transactions on Neural Networks} 9, 987--1000 (1998). \href{https://doi.org/10.1109/72.712178}{doi:10.1109/72.712178}.
\bibitem{ns} OpenAI (attributed author in supplied manuscript). \emph{Finite Time Blowup for Navier--Stokes}. User-supplied file \texttt{navier-stokes-c,d.pdf}; Section 5 and Appendix B. Used as methodological motivation; publication status and the complete proof have not been independently verified here.
\end{thebibliography}
\end{document}
